\section{实战案例：落地页文案 Skill 优化}

\subsection{优化结果}

作者在自己的落地页文案 skill 上运行 autoresearch，结果如下：

$$\text{通过率：} 56\% \longrightarrow 92\%$$

共 4 轮改动（rounds of changes），其中 3 个保留，1 个撤销。

\subsection{Agent 具体做了哪些改动}

\begin{enumerate}
    \item \textbf{针对最高频失败项加了明确规则}：``Your headline must include a specific number or result. Never use vague promises like `Transform Your Business.'\,'' ——直接堵住最常见的失败模式。

    \item \textbf{加了禁用词列表（banned buzzwords list）}：明确列出禁止使用的词汇——revolutionary, cutting-edge, synergy, next-level, game-changing, leverage, unlock, transform。

    \item \textbf{加了一个 worked example}：在 skill prompt 中加入一段实际的高质量落地页示例，标注了痛点开场白（pain point opener）和 CTA 的位置，让 skill 能\textbf{直接看到``好''的样子}，而不是靠猜。

    \item \textbf{尝试更严格的字数限制（被撤销）}：文案变得太单薄（too thin），CTA 质量也下降了。系统检测到了这种``单独看似改进、但实际损害整体输出''的改动，自动撤销。
\end{enumerate}

\begin{knowledgebox}{改动被撤销的智慧}
第 4 轮改动展示了 autoresearch 的一个关键能力：它不是盲目叠加约束，而是能识别\textbf{局部改进但全局退步}的情况。更严的字数限制让``字数在 150 词以内''这项检查更容易通过，但牺牲了 CTA 质量和内容丰富度。系统权衡了整体得分后选择撤销。
\end{knowledgebox}

\subsection{最终产出物}

运行结束后，你会得到以下 4 样东西：

\begin{enumerate}
    \item \textbf{改进后的 skill}——保存为单独文件，原版完好无损，随时可回退。
    \item \textbf{结果日志（results log）}——记录每一轮的得分。
    \item \textbf{Changelog}——解释每次改动的内容、agent 为什么尝试这个改动、以及改动是否有效。
    \item \textbf{原始 skill 的备份}——以防你需要回到起点。
\end{enumerate}

\begin{importantbox}{Changelog 是最有价值的产出}
那份 changelog 可能是整个过程中最值钱的东西。它是一份完整的记录：对于这个特定的 skill，什么有效、什么无效。

更重要的是，当更强的模型问世时，你把这份 changelog 交给新模型，它就能从上一个 agent 停下的地方\textbf{接着优化}。Changelog 让 skill 的改进经验变成了可传承的资产。
\end{importantbox}

\subsection{本章小结}

实战结果证明了 autoresearch 的有效性：4 轮自动迭代将通过率从 56\% 提升到 92\%。系统能自动识别有害改动并撤销。最终产出不仅是优化后的 skill，还包括完整的 changelog，为后续持续优化和跨模型传承提供基础。

